%% -*- coding: utf-8 -*-
\documentclass{beamer}

%% -*- coding: utf-8 -*-
\usetheme{Boadilla} % default
\useoutertheme{infolines}
\setbeamertemplate{navigation symbols}{} 

\usepackage{etex}
\usepackage{alltt}
\usepackage{pifont}
\usepackage{color}
\usepackage[utf8]{inputenc}
%\usepackage{german}
\usepackage{listings}
\definecolor{keywordcolor}{rgb}{0.7,0.1,0.1}
\definecolor{tacticcolor}{rgb}{0.0,0.1,0.6}
\definecolor{commentcolor}{rgb}{0.4,0.4,0.4}
\definecolor{symbolcolor}{rgb}{0.0,0.1,0.6}
\definecolor{sortcolor}{rgb}{0.1,0.5,0.1}
\definecolor{attributecolor}{rgb}{0.7,0.1,0.1}

\def\lstlanguagefiles{lstlean.tex}
\lstset{language=Lean}
\lstset{sensitive=true}
\usepackage{hyperref}
\hypersetup{colorlinks=true}
\usepackage[final]{pdfpages}
\usepackage{url}
\usepackage{arydshln} % dashed lines
\usepackage{tikz}
\usepackage{mathpartir}

\DeclareUnicodeCharacter{3BB}{\ensuremath{\lambda}}


\newcommand\cmark{\ding{51}}
\newcommand\xmark{\ding{55}}

\newcommand{\nat}{\mathbf{N}}

\usepackage[all]{xy}

%% new arrow tip for xy
\newdir{|>}{!/4.5pt/@{|}*:(1,-.2)@^{>}*:(1,+.2)@_{>}}

\newcommand\cid[1]{\textup{\textbf{#1}}} % class names
\newcommand\kw[1]{\textup{\textbf{#1}}}  % key words
\newcommand\tid[1]{\textup{\textsf{#1}}} % type names
\newcommand\vid[1]{\textup{\texttt{#1}}} % value names
\newcommand\Mid[1]{\textup{\texttt{#1}}} % method names

\newcommand\TODO[1][]{{\color{red}{\textbf{TODO: #1}}}}

\newcommand\String[1]{\texttt{\dq{}#1\dq{}}}

\newcommand\ClassHead[1]{%
  \ensuremath{\begin{array}{|l|}
      \hline
      \cid{#1}
      \\\hline
    \end{array}}}
\newcommand\AbstractClass[2]{%
  \ensuremath{\begin{array}{|l|}
      \hline
      \cid{\textit{#1}}
      \\\hline
      #2
      \hline
    \end{array}}}
\newcommand\Class[2]{%
  \ensuremath{\begin{array}{|l|}
      \hline
      \cid{#1}
      \\\hline
      #2
      \hline
    \end{array}}}
\newcommand\Attribute[3][black]{\textcolor{#1}{\Param{#2}{#3}}\\}
\newcommand\Methods{\hline}
\newcommand\MethodSig[3]{\Mid{#2} (#3): \,\tid{#1}\\}
\newcommand\CtorSig[2]{\Mid{#1} (#2)\\}
\newcommand\AbstractMethodSig[3]{\Mid{\textit{#2}} (#3): \,\tid{#1}\\}
\newcommand\Param[2]{\vid{#2}:~\tid{#1}}

\lstset{%
  frame=single,
  xleftmargin=2pt,
  stepnumber=1,
  numbers=left,
  numbersep=5pt,
  numberstyle=\ttfamily\tiny\color[gray]{0.3},
  belowcaptionskip=\bigskipamount,
  captionpos=b,
  escapeinside={*'}{'*},
  % language=java,
  tabsize=2,
  emphstyle={\bf},
  commentstyle=\mdseries\it,
  stringstyle=\mdseries\rmfamily,
  showspaces=false,
  showtabs=false,
  keywordstyle=\bfseries,
  columns=fullflexible,
  basicstyle=\footnotesize\CodeFont,
  showstringspaces=false,
  morecomment=[l]\%,
  rangeprefix=////,
  includerangemarker=false,
}

\newcommand\CodeFont{\sffamily}

\definecolor{lightred}{rgb}{0.8,0,0}
\definecolor{darkgreen}{rgb}{0,0.5,0}
\definecolor{darkblue}{rgb}{0,0,0.5}

\newcommand\highlight[1]{\textcolor{blue}{\emph{#1}}}
\newcommand\GenClass[2]{\cid{#1}\texttt{<}\cid{#2}\texttt{>}}

\newcommand\Colored[3]{\alt<#1>{\textcolor{#2}{#3}}{#3}}

\newcommand\nt[1]{\ensuremath{\langle#1\rangle}}

\newcommand{\free}{\operatorname{free}}
\newcommand{\bound}{\operatorname{bound}}
\newcommand{\var}{\operatorname{var}}
\newcommand\VSPBLS{\vspace{-\baselineskip}}

\newcommand\IF{\textit{IF}}
\newcommand\TRUE{\textit{TRUE}}
\newcommand\FALSE{\textit{FALSE}}

\newcommand\IFZ{\textit{IF0}}
\newcommand\ZERO{\textit{ZERO}}
\newcommand\SUCC{\textit{SUCC}}
\newcommand\ADD{\textit{ADD}}
\newcommand\SUB{\textit{SUB}}
\newcommand\MULT{\textit{MULT}}
\newcommand\DIV{\textit{DIV}}

\newcommand\PAIR{\textit{PAIR}}
\newcommand\FST{\textit{FST}}
\newcommand\SND{\textit{SND}}

\newcommand\CASE{\textit{CASE}}
\newcommand\LEFT{\textit{LEFT}}
\newcommand\RIGHT{\textit{RIGHT}}

\newcommand\Encode[1]{\lceil#1\rceil}
\newcommand\Reduce{\stackrel\ast\rightarrow_\beta}

\newcommand\Nat{\textit{Nat}}
\newcommand\Bool{\textit{Bool}}
\newcommand\Pair{\textit{Pair}}
\newcommand\Tfun[1]{#1\to}

\newcommand\Tenv{A}
\newcommand\Lam[1]{\lambda#1.}
\newcommand\App[1]{#1\,}
\newcommand\Succ{\textit{SUCC}\,}
\newcommand\Let[2]{\textit{let}\,#1=#2\,\textit{in}\,}

\newcommand\calE{\mathcal{E}}
\newcommand\calU{\mathcal{U}}
\newcommand\calP{\mathcal{P}}
\newcommand\calW{\mathcal{W}}

\newcommand\GEN{\textit{gen}}
\newcommand\EFV[1]{\textit{fv} (#1)}
\newcommand\Dom[1]{\textit{dom} (#1)}

%%% Local Variables: 
%%% mode: latex
%%% TeX-master: nil
%%% End: 

%%% frontmatter
%% -*- coding: utf-8 -*-

\title{Functional Programming / Lean}
\subtitle{Introduction}

\author[Peter Thiemann]{Prof. Dr. Peter Thiemann}
\institute[Univ. Freiburg]{Albert-Ludwigs-Universität Freiburg, Germany}
\date{WS 2026/27}


\subtitle{IO}
\usepackage{tikz}


\begin{document}

\begin{frame}
  \titlepage
\end{frame}
%----------------------------------------------------------------------
\begin{frame}[fragile]
  \frametitle{Referential transparency and substitutivity}
  \begin{block}<+->{Remember the first class?}
    \begin{itemize}
    \item Every variable and expression has just one value\\
      \textbf{referential transparency}
    \item Every variable can be replaced by its definition\\
      \textbf{substitutivity}
    \end{itemize}
  \end{block}
  \begin{block}<+->{Enables reasoning}
\begin{lstlisting}
-- sequence of function calls does not matter
f () + g () == g () + f ()
-- number of function calls does not matter
f () + f () == 2 * f ()
\end{lstlisting}
  \end{block}
\end{frame}
%----------------------------------------------------------------------
\begin{frame}[fragile]
  \frametitle{How does IO fit in?}
  \begin{alertblock}<+->{Bad example}
    Suppose we had 
\begin{lstlisting}
input : Unit -> String
\end{lstlisting}
    \begin{itemize}
    \item<+-> Consider
\begin{lstlisting}
let x := input ()
x ++ x
\end{lstlisting}
    \item<+-> Expect to read one input and use it twice
    \item<+-> By substitutivity, this expression must behave like
\begin{lstlisting}
input () ++ input ()
\end{lstlisting}
      which reads two inputs!
    \item<+-> VERY WRONG!!!
  \end{itemize}
  \end{alertblock}
\end{frame}
%----------------------------------------------------------------------
\begin{frame}[fragile]
  \frametitle{The dilemma}
  \begin{alertblock}<+->{Lean is a pure language, but IO is a side effect}
  \end{alertblock}
  \begin{block}<+->{A contradiction?}
  \end{block}
  \begin{block}<+->{No!}
    \begin{itemize}
    \item Instead of performing IO operations directly, there is an
      abstract type of \textbf{IO instructions}, which get executed
      lazily by the operating system
    \item Some instructions (e.g., read from a file) return values, so the abstract IO type is parameterized over their type
    \item Keep in mind: instructions are just values like any other
  \end{itemize}
  \end{block}
\end{frame}
%----------------------------------------------------------------------
\begin{frame}[fragile]
  \frametitle{Lean IO}

\begin{block}<+->{The main function}
  Top-level result of a program is an IO ``instruction''.
\begin{lstlisting}
def main : IO Unit := sorry
\end{lstlisting}
  \begin{itemize}
  \item an instruction describes the \textbf{effect} of the program
  \item effect $=$ IO action, imperative state change, \dots
  \end{itemize}
\end{block}
\begin{block}<+->{An instruction that returns a result}
\begin{lstlisting}
IO.getStdout : IO IO.FS.Stream
\end{lstlisting}
\end{block}
\end{frame}
%----------------------------------------------------------------------
\begin{frame}[fragile]
  \frametitle{Kinds of instructions}
  \begin{block}<+->{Primitive instructions}
\begin{lstlisting}
open IO.FS.Stream
putStrLn : IO.FS.Stream → String → IO Unit
getLine  : (self : IO.FS.Stream) → IO String
IO.FS.readFile : (fname : System.FilePath) → IO String
IO.FS.writeFile : (fname : System.FilePath) → (content : String) → IO Unit
\end{lstlisting}
and many more (\lstinline|String| converts to \lstinline|System.FilePath|)
  \end{block}
  \begin{block}<+->{No op instruction}
\begin{lstlisting}
pure : α → IO α
\end{lstlisting}
    The IO instruction \lstinline{pure 42} performs no IO, but yields the value
    42.  \lstinline{return} is mostly synonymous to \lstinline{pure}.
  \end{block}
\end{frame}
%----------------------------------------------------------------------
\begin{frame}[fragile]
  \frametitle{Combining two instructions}
  \begin{block}<+->{The bind operator \texttt{>>=}}
    Intuition:  next instruction may depend on the output of the previous one
\begin{lstlisting}
bind : IO α → (α → IO β) → IO β
\end{lstlisting}
    The instruction \lstinline{bind m f} ( or infix \lstinline{m >>= f})
    \begin{itemize}
    \item executes \texttt{m : IO a} first
    \item gets its result \texttt{x : a}
    \item applies \texttt{f : a -> IO b} to the result
    \item to obtain an instruction \texttt{f x : IO b} that returns a \texttt{b}
    \item and executes this instruction to return a \texttt{b}
    \end{itemize}
  \end{block}
  \begin{exampleblock}<+->{Example}
\begin{lstlisting}
def readFiles (f1 f2 : System.FilePath) : IO String :=
  readFile f1 >>= fun s1 => readFile f2 >>= fun s2 =>
  pure (s1 ++ s2)
\end{lstlisting}
  \end{exampleblock}
\end{frame}
%----------------------------------------------------------------------
\begin{frame}[fragile]
  \frametitle{More convenient: \texttt{do} notation}
\begin{lstlisting}
def readFiles' (f1 f2 : System.FilePath) : IO String := do
  let s1 ← readFile f1
  let s2 ← readFile f2
  pure (s1 ++ s2)

def copyFile (src dst : System.FilePath) : IO Unit :=
  sorry

def doTwice (action : IO Unit) : IO Unit :=
  sorry

def doNot (action : IO Unit) : IO Unit :=
  sorry
\end{lstlisting}
\end{frame}
%----------------------------------------------------------------------
\begin{frame}
  \frametitle{Translating \texttt{do} notation into \texttt{>>=} operations}
  \begin{itemize}
  \item \texttt{do \emph{lastinstruction}}
    \\ $\longrightarrow$
    \\ \texttt{\emph{lastinstruction}}
    \\[2ex]
  \item \texttt{do \{ let \emph{x} <- \emph{action1}; \emph{instructions} \}}
    \\ $\longrightarrow$
    \\ \texttt{\emph{action1} >>= fun \emph{x} -> \texttt{do \{ \emph{instructions} \}}}
    \\[2ex]
  \item \texttt{do \{ \emph{action1}; \emph{instructions} \}}
    \\ $\longrightarrow$
    \\ \texttt{\emph{action1} >> \texttt{do \{ \emph{instructions} \}}}
    \\[2ex]
  \item \texttt{do \{ let \emph{binding}; \emph{instructions} \}}
    \\ $\longrightarrow$
    \\ \texttt{let \emph{binding} in do \{  \emph{instructions} \}}
  \end{itemize}
\end{frame}
%----------------------------------------------------------------------
\begin{frame}[fragile]
  \frametitle{Instructions vs functions}
  \begin{block}<+->{Functions}
    produce the same result each time they are called
  \end{block}
  \begin{block}<+->{Instructions}
    \begin{itemize}
    \item are interpreted each time they are executed,
    \item the result depends on the context
    \item may be different each time
    \end{itemize}
  \end{block}
\end{frame}
%----------------------------------------------------------------------
\begin{frame}
  \frametitle{Underlying concept: \textbf{Monad}}

  \begin{block}<+->{What's a monad? (first approximation)}
    \begin{itemize}
    \item abstract datatype for instructions that produce values
    \item built-in combination \texttt{>>=}
    \item abstracts over different interpretations (computations)
    \end{itemize}
  \end{block}
  \begin{alertblock}<+->{IO is a special case of a monad}
    \begin{itemize}
    \item \lstinline{Monad} is a type class that abstracts an API for monads
    \item \lstinline{IO} is one useful instance of \lstinline{Monad}
    \item built into Lean
    \item but there's more to the concept
    \item many more instances to come!
    \end{itemize}
  \end{alertblock}
\end{frame}
% %----------------------------------------------------------------------
% \begin{frame}[fragile]
%   \frametitle{Hands-on task}
%   Define a function
% \begin{lstlisting}
% sortFile : FilePath -> FilePath -> IO Unit

% -- sortFile inFile outFile 
% -- reads inFile, sorts its lines, and writes the result to outFile

% -- recall
% -- sort : Ord a => [a] -> [a]
% -- lines : String -> [String]
% -- unlines : [String] -> String
% \end{lstlisting}
% \end{frame}
% %----------------------------------------------------------------------
% \begin{frame}[fragile]
%   \frametitle{Utilities}
% \begin{lstlisting}
% sequence : [IO a] -> IO [a]
% sequence_ : [IO a] -> IO Unit
% \end{lstlisting}
% \end{frame}
% %----------------------------------------------------------------------
% \begin{frame}[fragile]
%   \frametitle{Another hands-on task}
%   Define a function
% \begin{lstlisting}
% printTable : [String] -> IO Unit

% {-
% printTable ["New York", "Rio", "Tokio"]
% outputs
% 1: New York
% 2: Rio
% 3: Tokio
% -} 
% \end{lstlisting}
% \end{frame}
% %----------------------------------------------------------------------
%----------------------------------------------------------------------
\begin{frame}
  \frametitle{Wrapup}
  \begin{alertblock}{First encounter with monads}
  \begin{itemize}
  \item A monad is an abstract data type of instructions returning results.
  \item The next instruction can depend on previous results.
  \item Instructions are just values.
  \item Lean's IO operations are instructions of the \texttt{IO} monad.
  \end{itemize}
  \end{alertblock}
\end{frame}

% \begin{frame}
%   \frametitle{Break Time --- Questions?}
%   \begin{center}
%     \tikz{\node[scale=15] at (0,0){?};}
%   \end{center}
% \end{frame}


\end{document}

%%% Local Variables: 
%%% mode: latex
%%% TeX-master: t
%%% End: 
